\section{Features and Capabilities}
\label{sec:features}
\begin{itemize}
\item \textbf{Dual-process chat.} Every reply shows which path produced it,
the assurance level, the thinking budget used and the cost; the user can move
a single ``speed--care'' control, which sets $\lambda$
(Proposition~\ref{prop:frontier}).
\item \textbf{Verified answers on code, mathematics and structured output},
with precision bounded by the verifier (Theorem~\ref{thm:precision}), and
evidence-fed retries (Theorem~\ref{thm:gatedloop}).
\item \textbf{Calibrated abstention.} When the budget is exhausted the system
says so, returns the best candidate with an uncertainty notice, and never
upgrades consensus to verification (Theorem~\ref{thm:plurality}).
\item \textbf{Honest self-description.} A per-domain capability table from
audited traces, with unmeasured domains marked as such
(Proposition~\ref{prop:audit}).
\item \textbf{Performance console.} Load tests against the gateway with
latency quantiles, throughput, escalation rate and cost per request, run from
the browser.
\item \textbf{Evaluation dashboard.} Benchmark runs with paired comparisons
against the fast-only and slow-only baselines, confidence intervals, routing
accuracy and calibration plots, and the go/no-go status of each roadmap gate.
\item \textbf{Gradual distillation.} Verified slow-path traces become
training data for the fast path under a statistically sized promotion gate
with human sign-off (Proposition~\ref{prop:distill}).
\item \textbf{Backend portability.} The same gateway runs against
llama-server on a laptop and vLLM on a server, CPU or GPU.
\end{itemize}
Honest limits: non-verifiable domains never exceed the \texttt{consensus}
assurance level; the fast and slow paths share a model family, so their errors
are correlated and the collision term of Proposition~\ref{prop:collision} is
measured, not assumed small; on the reference machine the slow path is not
interactive.
